\documentclass[11pt]{article}
\usepackage[a4paper,margin=1in]{geometry}
\usepackage{amsmath,amssymb,amsthm,mathtools,mathrsfs}
\usepackage{array,booktabs,longtable}
\usepackage[hidelinks]{hyperref}
\usepackage[T1]{fontenc}
\usepackage{lmodern}

\newtheorem{theorem}{Theorem}[section]
\newtheorem{lemma}[theorem]{Lemma}
\newtheorem{proposition}[theorem]{Proposition}
\newtheorem{corollary}[theorem]{Corollary}
\theoremstyle{definition}
\newtheorem{definition}[theorem]{Definition}
\theoremstyle{remark}
\newtheorem{remark}[theorem]{Remark}

\newcommand{\C}{\mathbb C}
\newcommand{\D}{\mathbb D}
\newcommand{\T}{\mathbb T}
\newcommand{\spec}{\operatorname{spec}}
\newcommand{\adj}{\operatorname{adj}}
\newcommand{\conv}{\operatorname{conv}}
\newcommand{\Rea}{\operatorname{Re}}
\newcommand{\Ima}{\operatorname{Im}}
\newcommand{\Span}{\operatorname{span}}
\newcommand{\dist}{\operatorname{dist}}
\newcommand{\Int}{\operatorname{int}}
\newcommand{\ol}[1]{\overline{#1}}

\title{Rigidity of the Extremal Crouzeix Ratio}
\author{}
\date{}

\begin{document}
\maketitle

\begin{abstract}
We prove that the assertion in the problem is true.  The proof includes the
sharp estimate with constant \(2\), so it does not assume the Crouzeix
conjecture.  Equality is first converted, without an attainment assumption on
polynomials, into a boundary-kernel identity for a finite Blaschke extremizer.
That identity makes the extremizer rational.  A positive extremal state then
shows that its entire unit lemniscate, not merely one arc, is the boundary of
the numerical range.  Finally, the dual of this boundary curve is a factor of
a Hermitian characteristic polynomial.  Hyperbolicity and strict convexity
force the dual to be a conic, and the two circular points at infinity force the
conic to be a circle.
\end{abstract}

\section{Problem and notation}

For \(A\in M_N(\C)\), set
\[
 W(A)=\{u^*Au:u\in\C^N,\ \|u\|_2=1\}.
\]
For a polynomial \(p\), let \(\|p(A)\|\) be its Euclidean operator norm and put
\[
 \psi(A)=
 \sup_{\substack{p\in\C[z]\\ \max_{z\in W(A)}|p(z)|>0}}
 \frac{\|p(A)\|}{\max_{z\in W(A)}|p(z)|}.
\]
We prove the following statement, including the strict positivity of the
radius in its conclusion.

\begin{theorem}[Main theorem]\label{thm:main}
For every \(N\geq1\) and \(A\in M_N(\C)\),
\[
 \psi(A)=2\quad\Longrightarrow\quad
 W(A)=\{z\in\C:|z-\gamma|\leq\rho\}
\]
for some \(\gamma\in\C\) and some \(\rho>0\).
\end{theorem}

Here is a notation table.  Scalar matrix coefficients are always written as
\(x^*Yy\), avoiding any dependence on an inner-product convention.

\begin{center}
\begin{tabular}{@{}p{0.25\textwidth}p{0.68\textwidth}@{}}
\toprule
Symbol & Meaning\\
\midrule
\(K\) & the compact convex set \(W(A)\)\\
\(\Omega\) & the interior of \(K\), when nonempty\\
\(R_z\) & the resolvent \((zI-A)^{-1}\)\\
\(H_A=\Rea A\), \(G_A=\Ima A\) & \((A+A^*)/2\), respectively \((A-A^*)/(2i)\)\\
\(h_A(\theta)\) & \(\lambda_{\max}(\Rea(e^{-i\theta}A))\)\\
\(A(K)\) & functions continuous on \(K\) and holomorphic on \(\Omega\)\\
\(\psi_\Omega(A_0)\) & bounded-holomorphic calculus norm of a matrix \(A_0\) on \(\Omega\)\\
\(\Delta_\sigma,\mathcal P(\sigma)\) & support defect and positive double-layer density at \(\sigma\in\partial K\)\\
\(T,x,y\) & the extremal matrix \(T=f(A)\) and its equality vectors\\
\(a(z),c(z)\) & the diagonal and off-diagonal resolvent entries in (5.2)\\
\(U,V\) & coprime numerator and denominator of the reduced rational extremizer\\
\(\mathscr F,\mathscr C,\mathscr G\) & homogenized lemniscate, its boundary factor, and the dual factor\\
\(S^\#(w)\) & \(\overline{S(\overline w)}\) for a polynomial \(S\)\\
\bottomrule
\end{tabular}
\end{center}

\section{Foundational facts}

\begin{lemma}[Numerical-range geometry]\label{lem:foundation}
The following facts hold.
\begin{enumerate}
\item \(W(A)\) is nonempty, compact, and convex.
\item If \(W(A)\) is a singleton or is contained in a line, then \(A\) is
normal and \(\psi(A)=1\).
\item For \(a\neq0\), \(b\in\C\), and unitary \(\mathcal U\),
\[
 W(aA+bI)=aW(A)+b,\qquad \psi(aA+bI)=\psi(A),\qquad
 \psi(\mathcal U^*A\mathcal U)=\psi(A).
\]
Also \(W(A^*)=\overline{W(A)}\) and \(\psi(A^*)=\psi(A)\).
\item For a finite direct sum,
\[
 W\left(\bigoplus_jA_j\right)=\conv\Bigl(\bigcup_jW(A_j)\Bigr),
 \qquad
 \psi\left(\bigoplus_jA_j\right)\leq\max_j\psi(A_j).
\]
\item The support function of \(W(A)\) is
\[
 h_A(\theta)=\lambda_{\max}\!\left(\Rea(e^{-i\theta}A)\right).
\]
\end{enumerate}
\end{lemma}

\begin{proof}
The unit sphere is compact and \(u\mapsto u^*Au\) is continuous, proving
nonemptiness and compactness.  To prove convexity, take two points of \(W(A)\)
and compress \(A\) to the span of corresponding unit vectors.  It is enough to
know that the numerical range of a \(2\times2\) matrix is convex.  After unitary
triangularization such a matrix has the form
\(
 \begin{psmallmatrix}\lambda_1&c\\0&\lambda_2\end{psmallmatrix}
\).
For \(u=(\cos t,e^{i\varphi}\sin t)\),
\[
 u^*Au=\lambda_1\cos^2t+\lambda_2\sin^2t
       +c e^{i\varphi}\sin t\cos t.
\]
These points fill the possibly degenerate ellipse with foci
\(\lambda_1,\lambda_2\), hence a convex set.  The two-dimensional compression
is contained in \(W(A)\), so the segment joining the original two points is
contained in \(W(A)\).

If \(W(A)=\{\gamma\}\), then \(x^*(A-\gamma I)x=0\) for all \(x\).  The complex
polarization identity makes every matrix coefficient of \(A-\gamma I\) zero.
If \(W(A)\) lies in a line, a translation and a nonzero complex rotation make
all quadratic forms of \(A\) real.  The skew-Hermitian part then has zero
quadratic form and vanishes by polarization, so the transformed matrix is
Hermitian.  A normal matrix satisfies
\[
 \|p(A)\|=\max_{\lambda\in\spec(A)}|p(\lambda)|
 \leq\max_{z\in W(A)}|p(z)|.
\]
The nonzero constant polynomials give the reverse inequality \(1\).  This also
settles the positivity convention: a nonzero constant has positive
denominator, whereas polynomials with zero denominator are not used.

The affine and unitary formulas follow by changing variables in polynomials.
For adjoints use \(\widetilde p(z)=\overline{p(\overline z)}\), for which
\(\widetilde p(A^*)=p(A)^*\).  For a direct sum, a unit vector decomposes into block
components and its numerical value is a convex combination of block numerical
values.  Also \(p(\oplus A_j)=\oplus p(A_j)\); the supremum norm on the convex
hull dominates that on every block.  Finally,
\[
 \max_{z\in W(A)}\Rea(e^{-i\theta}z)
 =\max_{\|u\|=1}u^*\Rea(e^{-i\theta}A)u,
\]
which is the largest eigenvalue of the displayed Hermitian matrix.
\end{proof}
\section{A self-contained sharp bound and its equality case}

We next prove the sharp constant \(2\).  It is therefore not an assumption in
the rigidity argument.

\begin{lemma}[An abstract dilation estimate]\label{lem:dilation}
Let \(T\) act on a finite-dimensional Hilbert space.  Suppose that \(\mathcal M\) is a
contraction on another Hilbert space, \(\mathcal V\) is an isometry, and
\[
 E_n=2\mathcal V^*\mathcal M^{*n}\mathcal V-T^{*n}\qquad(n\geq1)
\]
are uniformly bounded and commute with \(T\).  Then \(\|T\|\leq2\).

If additionally \(\|T\|=2\), \(r(T)<1\), \(T^*Tx=4x\), \(\|x\|=1\), and
\(y=Tx/2\), then
\[
 T^*x=0,\quad T^*y=2x,\quad
 \mathcal M\mathcal Vx=\mathcal Vy,\quad
 \mathcal M^*\mathcal Vy=\mathcal Vx.
\]
\end{lemma}

\begin{proof}
Put \(\kappa=\|T\|\).  There is nothing to prove if \(\kappa\leq1\).  Choose a
unit vector \(x\) with \(T^*Tx=\kappa^2x\), and set
\[
 m_n=\Rea\bigl(x^*E_nT^nx\bigr),\qquad
 \omega=\mathcal M^*\mathcal VTx-\kappa\mathcal Vx.
\]
Indeed
\(\|T^{*n}\|\leq2+\sup_k\|E_k\|\) by the defining identity, so the powers of
\(T\), and hence \(m_n\), are uniformly bounded.  Commutativity and direct
expansion give
\begin{align*}
 \kappa m_n-m_{n+1}
 &=\kappa(\kappa-1)\|T^{*n}x\|^2
   -2\Rea\bigl((\mathcal V^*\mathcal M^{*n}\omega)^*T^{*n}x\bigr)\\
 &=\kappa(\kappa-1)
   \left\|T^{*n}x-\frac{\mathcal V^*\mathcal M^{*n}\omega}
   {\kappa(\kappa-1)}\right\|^2
   -\frac{\|\mathcal V^*\mathcal M^{*n}\omega\|^2}{\kappa(\kappa-1)}\\
 &\geq-\frac{\|\omega\|^2}{\kappa(\kappa-1)}.
\end{align*}
Telescoping, with the boundedness of \(m_n\), yields
\[
 \kappa m_1
 =\sum_{n\geq1}\kappa^{-n+1}(\kappa m_n-m_{n+1})
 \geq-\frac{\|\omega\|^2}{(\kappa-1)^2}.               \tag{3.1}
\]
On the other hand, using that \(\mathcal V\) is isometric and \(\mathcal M\) contractive,
\begin{align*}
 \|\omega\|^2
 &\leq2\kappa^2
 -2\kappa\Rea\bigl(x^*\mathcal V^*\mathcal M^*\mathcal VTx\bigr)\\
 &=2\kappa^2-\kappa m_1-\kappa^3.                     \tag{3.2}
\end{align*}
Combining (3.1) and (3.2) gives
\[
 \|\omega\|^2\left(1-\frac1{(\kappa-1)^2}\right)
 \leq\kappa^2(2-\kappa),
\]
which is impossible for \(\kappa>2\).

Now let \(\kappa=2\).  Equations (3.1)--(3.2) become
\(2m_1\geq-\|\omega\|^2\) and \(\|\omega\|^2\leq-2m_1\), so equality holds in both.
For each \(n\), the slack in the estimate preceding (3.1) is
\[
 2\left\|T^{*n}x-\tfrac12\mathcal V^*\mathcal M^{*n}\omega\right\|^2
 +\tfrac12\bigl(\|\omega\|^2-\|\mathcal V^*\mathcal M^{*n}\omega\|^2\bigr).
\]
All these nonnegative slacks have a positive weighted sum equal to zero.
Thus each vanishes.  Since \(r(T)<1\), \(T^{*n}x\to0\); the equality of norms
then forces \(\omega=0\), and the first term with \(n=1\) gives \(T^*x=0\).
With \(y=Tx/2\), \(\mathcal M^*\mathcal Vy=\mathcal Vx\).  Both
\(\mathcal Vx\) and \(\mathcal Vy\) are unit vectors, so equality in
Cauchy--Schwarz for \(\mathcal M\) gives \(\mathcal M\mathcal Vx=\mathcal Vy\).
Finally \(T^*y=T^*Tx/2=2x\).
\end{proof}

\begin{proposition}[The sharp numerical-range estimate]\label{prop:sharp}
For every finite matrix \(A\) and every polynomial \(p\),
\[
 \|p(A)\|\leq2\max_{z\in W(A)}|p(z)|.                 \tag{3.3}
\]
More generally, if \(K\) is a compact convex body,
\(W(A)\subset K\), \(\spec(A)\subset\Int K\), and \(f\in A(K)\), then
\[
 \|f(A)\|\leq2\|f\|_K.                               \tag{3.4}
\]
The same holds for \(f\in H^\infty(\Int K)\).
\end{proposition}

\begin{proof}
Normalize \(\|f\|_K\leq1\).  Cauchy's boundary formula below also holds when
\(f\) is only in \(A(K)\): for \(z_0\in\Int K\), the functions
\(f(z_0+r(z-z_0))\), \(r<1\), are holomorphic on a neighborhood of \(K\) and
converge uniformly to \(f\); apply the usual formula to them and pass to the
limit.  Parametrize the rectifiable convex Jordan curve
\(\Gamma=\partial K\) counterclockwise by arclength.  Its outward unit normal
\(\mathfrak n(\sigma)\) exists almost everywhere.  Put
\(R_\sigma=(\sigma I-A)^{-1}\) and
\[
 \mathcal P(\sigma)=\frac1\pi\Rea\bigl(\mathfrak n(\sigma)R_\sigma\bigr).
\]
At a point with a normal, the supporting-half-plane property gives
\[
 \Delta_\sigma:=\Rea\bigl(\overline{\mathfrak n(\sigma)}(\sigma I-A)\bigr)\succeq0.
\]
Multiplication by the two resolvents shows
\[
 \Rea(\mathfrak nR_\sigma)=R_\sigma^*\Delta_\sigma R_\sigma\succeq0. \tag{3.5}
\]
Since \(d\sigma=i\mathfrak n(\sigma)\,ds\), Cauchy's formula and its adjoint give
\[
 \int_\Gamma \mathcal P(\sigma)\,ds=2I.              \tag{3.6}
\]
Consequently
\[
 (\mathcal Vx)(\sigma)=2^{-1/2}\mathcal P(\sigma)^{1/2}x
\]
defines an isometry into \(L^2(\Gamma;\C^N)\).  Let \(M_f\) be multiplication
by \(f\), and set \(T=f(A)\).  Splitting the two terms in \(\mathcal P\) and using Cauchy's
formula gives, for every \(n\geq1\),
\[
 E_n:=2\mathcal V^*M_f^{*n}\mathcal V-T^{*n}
 =\frac1{2\pi i}\int_\Gamma
   \overline{f(\sigma)^n}R_\sigma\,d\sigma.           \tag{3.7}
\]
The right side commutes with \(T\), because every resolvent does.  Also
\[
 \sup_n\|E_n\|\leq
 \frac{\operatorname{length}(\Gamma)}{2\pi}
 \max_{\sigma\in\Gamma}\|R_\sigma\|<\infty.
\]
Lemma~\ref{lem:dilation} proves (3.4).

For an \(H^\infty\) function, choose a Riemann map
\(\phi:\Int K\to\D\) and replace \(f\) by
\[
 f_r(z)=f\bigl(\phi^{-1}(r\phi(z))\bigr),\qquad r<1.
\]
These functions belong to \(A(K)\), retain norm at most \(1\), and their
spectral jets converge to those of \(f\); hence \(f_r(A)\to f(A)\).

Finally, for an arbitrary \(A\), apply (3.4) to the outer parallel convex body
\(K_\varepsilon=W(A)+\varepsilon\overline\D\), which contains \(W(A)\) in its
interior.  For a polynomial \(p\),
\(\max_{K_\varepsilon}|p|\to\max_{W(A)}|p|\).  Letting
\(\varepsilon\downarrow0\) proves (3.3), including all boundary-spectrum and
nonsmooth cases.
\end{proof}

\begin{proposition}[Equality data]\label{prop:equality-data}
Let \(K=W(A)\) have nonempty interior, let \(\spec(A)\subset\Omega:=\Int K\),
and let \(f\in A(K)\) be nonconstant with
\[
 \|f\|_K=1,\qquad |f|=1\ \hbox{on }\partial K,\qquad \|f(A)\|=2.
\]
Then there are orthonormal vectors \(x,y\) such that, with \(T=f(A)\),
\[
 Tx=2y,\quad T^*x=0,\quad Ty=0,\quad T^*y=2x,         \tag{3.8}
\]
and, for almost every \(\sigma\in\partial K\),
\[
 \mathcal P(\sigma)^{1/2}\bigl(y-f(\sigma)x\bigr)=0.  \tag{3.9}
\]
Moreover, for every function \(h\) holomorphic near \(\spec(A)\) for which
the products below are defined,
\begin{align}
 x^*h(A)f(A)x&=0,& y^*h(A)f(A)x&=2x^*h(A)x,\tag{3.10}\\
 x^*h(A)y&=0,& y^*h(A)y&=x^*h(A)x.             \tag{3.11}
\end{align}
\end{proposition}

\begin{proof}
Apply the construction in Proposition~\ref{prop:sharp}.  The spectral mapping
theorem gives \(r(T)<1\), because the finitely many eigenvalues of \(A\) are
inside \(\Omega\) and a nonconstant inner function has modulus strictly less
than \(1\) there.  Choose \(x\) with \(T^*Tx=4x\), and put \(y=Tx/2\).
Lemma~\ref{lem:dilation} gives \(T^*x=0\), \(T^*y=2x\), and
\(M_f\mathcal Vx=\mathcal Vy\).  Here \(\|y\|=1\), and
\(2x^*y=x^*Tx=(T^*x)^*x=0\), so \(x,y\) are orthonormal.
The multiplication identity is exactly (3.9).

Apply the same equality argument to \(A^*\), the domain \(\overline K\), and
\(\widetilde f(z)=\overline{f(\overline z)}\).  Its matrix value is \(T^*\),
and \(y\) is a top right singular vector for \(T^*\).  The conclusion
\((T^*)^*y=0\) is \(Ty=0\).  Thus (3.8) holds.

Since \(T\) commutes with \(h(A)\), \(T^*x=0\) gives the first equality in
(3.10), and \(Tx=2y\) then gives the first equality in (3.11).  Also
\[
 x^*h(A)x=\tfrac12y^*Th(A)x
          =\tfrac12y^*h(A)Tx=y^*h(A)y,
\]
and multiplication by \(Tx=2y\) gives the second equality in (3.10).
\end{proof}


\begin{lemma}[Boundary eigenvalues reduce]\label{lem:boundary-eigen}
If \(Ax=\lambda x\), \(\|x\|=1\), and \(\lambda\in\partial W(A)\), then
\(A^*x=\overline\lambda x\).  Consequently
\[
 A\simeq A_{\partial}\oplus A_{\mathrm{in}},
\]
where \(A_{\partial}\) is diagonal and consists of all eigenvalues on
\(\partial W(A)\), while
\(\spec(A_{\mathrm{in}})\subset\Int W(A)\).
\end{lemma}

\begin{proof}
Choose a supporting direction \(n\), \(|n|=1\), at \(\lambda\).  Equality in
the Rayleigh quotient gives
\[
 \Rea(\overline nA)x=\Rea(\overline n\lambda)x.
\]
Substituting \(Ax=\lambda x\) gives \(A^*x=\overline\lambda x\).  Thus
\(\C x\) reduces \(A\).  Split it off and repeat.  Every eigenvalue belongs to
\(W(A)\), since an eigenvector realizes it, so the remaining spectrum is in
the interior.
\end{proof}

We shall use two elementary approximation facts.  Their proofs make all
passages between polynomials and holomorphic functions explicit.

\begin{lemma}[Polynomial approximation on a convex compact set]\label{lem:approx}
If \(K\subset\C\) is compact and convex with nonempty interior, then every
function in \(A(K)\) is a uniform limit on \(K\) of polynomials.
\end{lemma}

\begin{proof}
Fix \(z_0\in\Int K\).  For \(0<r<1\), put
\(f_r(z)=f(z_0+r(z-z_0))\).  Uniform continuity on \(K\) gives
\(f_r\to f\) uniformly, while \(f_r\) is holomorphic on a neighborhood of
\(K\).

For completeness, a function holomorphic near \(K\) can be approximated by
polynomials as follows.  Cauchy's formula on finitely many small polygonal
curves in that neighborhood, followed by Riemann sums, first approximates it
by a rational function whose poles lie in \(\C\setminus K\).  A pole can be
moved in sufficiently small steps along a polygonal path in
\(\C\setminus K\): the identity
\[
 \frac1{z-a}=\frac1{z-b}\,
 \frac1{1-(a-b)/(z-b)}
\]
and its geometric expansion are uniform on \(K\) whenever the step is shorter
than the distance to \(K\).  Subdivide the path and move every pole to a point
of arbitrarily large modulus.  At a large point \(b\),
\(1/(z-b)=-b^{-1}\sum_{j\geq0}(z/b)^j\) uniformly on \(K\), and truncation is
a polynomial.  The complement of a convex compact set is connected, so the
required paths exist.  This proves the neighborhood case and, on letting
\(r\uparrow1\), the lemma.
\end{proof}

\begin{lemma}[Conformal and finite-interpolation facts]\label{lem:complexfacts}
Let \(\Omega\) be a bounded convex domain.
\begin{enumerate}
\item There is a conformal bijection \(\phi:\Omega\to\D\), and it extends to a
homeomorphism \(\overline\Omega\to\overline\D\).
\item Given finitely many values and derivatives at points of \(\Omega\) that
are realized by a function of \(H^\infty(\Omega)\) of norm at most \(1\), the
same data are realized by \(\mathcal B\circ\phi\), where \(\mathcal B\) is a finite Blaschke
product.  If the least possible norm for nonconstant data is exactly \(1\),
the norm-one interpolant is unique and is of this form.
\end{enumerate}
\end{lemma}

\begin{proof}
We recall short proofs of the precise forms used here.  The Riemann mapping
theorem follows by maximizing the derivative at a fixed point in the normal
family of injective maps into \(\D\); a missing point permits a square-root
perturbation increasing that derivative, so the extremal map is onto.  A
bounded convex domain is a Jordan domain.  Images of nested crosscuts have
diameter tending to zero; otherwise two disjoint crosscuts would separate a
fixed interior compact set in incompatible ways.  Thus the Riemann map and
its inverse have unique continuous boundary limits.  The limits preserve
cyclic order, hence are mutually inverse.  This is the crosscut proof of the
Carath\'eodory extension specialized to a convex domain.

For interpolation, transfer the data by \(\phi\) to \(\D\), and list the nodes
with repetitions, a repetition encoding the next derivative.  If a Schur
function \(g\) has \(g(a)=b\), \(|b|<1\), Schwarz's lemma after two disk
automorphisms says that
\[
 \mathcal S_{a,b}g(z)=
 \frac{g(z)-b}{1-\overline b\,g(z)}
 \frac{1-\overline a z}{z-a}
\]
is again a Schur function.  At \(z=a\) the removable value records the next
derivative.  Repeating this operation consumes one interpolation condition at
a time.  After the last condition, choose the remaining Schur function to be
a unimodular constant and invert the transformations.  Each inversion adds a
disk automorphism and one Blaschke factor, so the result is a finite Blaschke
product satisfying all the original, including repeated-node, conditions.
If at some stage a parameter has modulus \(1\), Schwarz's lemma forces the
remaining function to be constant.

For the assertion about a least norm, divide the prescribed values and
derivatives by a common candidate bound \(c\), and run the same finite
algorithm.  Until a parameter reaches the unit circle, every prescribed
Schur parameter is a continuous function of \(c\); the data are feasible
exactly while these parameters have modulus at most \(1\).  If all of them
had modulus strictly less than \(1\) at the least feasible value, continuity
would leave them less than \(1\) after decreasing \(c\) slightly, a
contradiction.  Thus a first parameter has modulus \(1\).  Schwarz's lemma
then forces the remaining Schur function to be that unimodular constant, so
all subsequent choices are forced.  Inverting the preceding transformations
gives the unique interpolant, which is a finite Blaschke product.  This is
the finite Schur algorithm, including repeated nodes.
\end{proof}

For a bounded domain \(\Omega\) and a matrix \(A_0\) with
\(\spec(A_0)\subset\Omega\), define
\[
 \psi_\Omega(A_0)=\sup\{\|g(A_0)\|:g\in H^\infty(\Omega),\
                         \|g\|_\Omega\leq1\}.
\]
The value \(g(A_0)\) depends on only finitely many derivatives, those prescribed
by the Jordan blocks.  Montel's theorem therefore makes the supremum a
maximum.  Lemma~\ref{lem:complexfacts} replaces an extremizer by a finite
Blaschke product composed with a Riemann map, without changing \(g(A_0)\).

\begin{lemma}[Strict domain monotonicity]\label{lem:strict-domain}
Let \(\Omega_0\subsetneq\Omega_1\) be bounded convex domains and suppose
\(\spec(A_0)\subset\Omega_0\).  If \(\psi_{\Omega_1}(A_0)>1\), then
\[
 \psi_{\Omega_0}(A_0)>\psi_{\Omega_1}(A_0).
\]
\end{lemma}

\begin{proof}
Monotonicity in the displayed direction is immediate by restriction.  Suppose
equality held, with common value \(c>1\), and choose a norm-one extremizer \(g\)
on \(\Omega_1\).  Its norm on \(\Omega_0\) must be exactly \(1\), since otherwise
rescaling its restriction gives a value larger than \(c\).  Thus it is also an
extremizer on \(\Omega_0\).

For either domain, interpolate the finite jet of \(g\) that determines \(g(A_0)\)
with least possible norm.  A norm smaller than \(1\) would contradict
extremality after rescaling.  The uniqueness part of
Lemma~\ref{lem:complexfacts} says that \(g\) is a finite Blaschke product
relative to both domains.  Hence it has modulus \(1\) on each boundary.

Choose \(a\in\Omega_0\).  Along every ray from \(a\), convexity gives an exit
radius for each domain.  Since the domains differ, on some ray the first exit
occurs strictly earlier; hence there is
\(\xi\in\partial\Omega_0\cap\Omega_1\).  At this interior point of \(\Omega_1\),
\(|g(\xi)|=1\).  The maximum-modulus principle makes \(g\) constant, which
would give \(c=1\), a contradiction.
\end{proof}

\section{Reduction to an interior-spectrum equality pair}

\begin{lemma}[Boundary-spectrum induction]\label{lem:induction}
Assume Theorem~\ref{thm:main} has been proved in every size smaller than \(N\).
Let \(A\in M_N(\C)\), let \(K=W(A)\), and suppose \(\psi(A)=2\).  Then either
\(K\) is a nondegenerate closed disk, or
\[
 \spec(A)\subset\Int K.                               \tag{4.1}
\]
\end{lemma}

\begin{proof}
By Lemma~\ref{lem:foundation}, \(K\) has nonempty interior.  Choose
polynomials \(p_j\) such that
\[
 \max_K|p_j|=1,\qquad \|p_j(A)\|\longrightarrow2.      \tag{4.2}
\]
This is only an extremizing sequence; no polynomial maximum is being assumed.

Suppose \(A\) has a boundary eigenvalue.  Lemma~\ref{lem:boundary-eigen} gives
\(A\simeq A_\partial\oplus A_1\), where \(A_\partial\) is diagonal with
spectrum on \(\partial K\) and \(\spec(A_1)\subset\Int K\).  The
\(A_\partial\)-block of \(p_j(A)\) has norm at most
\(1\), so
\[
 \|p_j(A_1)\|\longrightarrow2.                        \tag{4.3}
\]
In particular \(A_1\) is nonempty.  Since \(W(A_1)\subset K\), (4.3) and
Proposition~\ref{prop:sharp} give \(\psi(A_1)=2\).  Its size is smaller than
\(N\), so induction says that
\[
 K_1:=W(A_1)
\]
is a closed disk of positive radius.

There is one more necessary peeling step.  Apply
Lemma~\ref{lem:boundary-eigen} to \(A_1\) relative to \(K_1\):
\[
 A_1\simeq A_{\partial,1}\oplus A_2,
\]
where \(A_{\partial,1}\) is a finite diagonal block on \(\partial K_1\) and
\(\spec(A_2)\subset\Int K_1\).  The same sequence (4.2) is bounded by \(1\) at
the eigenvalues of \(A_{\partial,1}\), all of which lie in \(K\), so
\(\|p_j(A_2)\|\to2\).  Thus \(A_2\) is nonempty.

Write \(\Lambda_1=\spec(A_{\partial,1})\).  Since
\[
 K_1=W(A_1)=\conv(\Lambda_1\cup W(A_2)),
\]
every extreme point of \(K_1\) belongs to the compact set
\(\Lambda_1\cup W(A_2)\).  Here is a direct verification of this last elementary
fact: if an extreme point of the convex hull of a compact set were not in the
set, a finite convex representation with at least two distinct points would
write it as a nontrivial segment point.  The circle \(\partial K_1\) has
infinitely many extreme points and \(\Lambda_1\) is finite.  Hence
\(\partial K_1\setminus\Lambda_1\subset W(A_2)\); closedness and convexity of
\(W(A_2)\) give
\[
 W(A_2)=K_1.                                           \tag{4.4}
\]

Let \(\Omega_0=\Int K_1\) and \(\Omega_1=\Int K\).  The polynomial sequence and
Proposition~\ref{prop:sharp} show
\[
 \psi_{\Omega_0}(A_2)=\psi_{\Omega_1}(A_2)=2.          \tag{4.5}
\]
Indeed the lower bounds follow from (4.2)--(4.4), while the upper bounds are
the relative estimate (3.4).  If \(K_1\subsetneq K\), then
\(\Omega_0\subsetneq\Omega_1\), and Lemma~\ref{lem:strict-domain} contradicts
(4.5).  Therefore \(K_1=K\), and \(K\) is the asserted disk.  If this conclusion
does not occur, \(A\) has no boundary eigenvalue, which is (4.1).
\end{proof}

\begin{lemma}[The holomorphic extremizer]\label{lem:extremizer}
Suppose \(K=W(A)\) has nonempty interior and
\(\spec(A)\subset\Omega=\Int K\).  If \(\psi(A)=2\), there is a nonconstant
\[
 f=\mathcal B\circ\phi\in A(K),
\]
where \(\phi:\Omega\to\D\) is conformal and \(\mathcal B\) is a finite
Blaschke product,
such that
\[
 \|f\|_K=1,\qquad |f|=1\text{ on }\partial K,\qquad
 \|f(A)\|=2.                                          \tag{4.6}
\]
\end{lemma}

\begin{proof}
Normalized polynomials witnessing \(\psi(A)=2\) show
\(\psi_\Omega(A)\geq2\), and Proposition~\ref{prop:sharp} gives the reverse
inequality.  Montel compactness, followed by convergence of the finitely many
spectral jets, gives a norm-one \(H^\infty(\Omega)\) maximizer.  Apply the
finite Schur algorithm in Lemma~\ref{lem:complexfacts} to the jet that
determines its matrix value.  It supplies \(\mathcal B\circ\phi\) with the same matrix
value.  The boundary extension of \(\phi\) puts this function in \(A(K)\) and
gives boundary modulus \(1\).  It cannot be constant because a constant matrix
has norm at most \(1\).  Finally Lemma~\ref{lem:approx} also shows directly
that this holomorphic maximum equals the original polynomial supremum.
\end{proof}

\section{From equality to a rational inner function}

We first locate an analytic boundary arc.  The following finite-dimensional
fact is included to avoid any hidden smoothness assumption.

\begin{lemma}[A curved numerical-range arc]\label{lem:arc}
If \(K=W(A)\) has interior, \(\spec(A)\subset\Int K\), and \(K\) is not a
polygon, then \(\partial K\) contains an open real-analytic arc on which every
supporting line exposes one point and the curvature is positive.
Moreover, under \(\spec(A)\subset\Int K\), \(K\) cannot be a polygon.
\end{lemma}

\begin{proof}
Put
\[
 \mathcal H(\theta)=\Rea(e^{-i\theta}A),\qquad
 h_A(\theta)=\lambda_{\max}(\mathcal H(\theta)).
\]
The roots of the characteristic polynomial of the analytic Hermitian pencil
\(\mathcal H(\theta)\) can locally be labeled by real-analytic functions.  One direct
proof is to split the spectral projections by small Cauchy circles and then
diagonalize the resulting analytic Hermitian blocks; repeated identical
branches are merged.  Two distinct analytic algebraic branches have only
finitely many crossings on the circle.  Thus the circle has a finite
partition into intervals on which \(h_A\) is analytic.

At a differentiability point, the exposed boundary point is
\[
 \sigma(\theta)=e^{i\theta}\bigl(h_A(\theta)+ih_A'(\theta)\bigr),\qquad
 \sigma'(\theta)=ie^{i\theta}\bigl(h_A(\theta)+h_A''(\theta)\bigr). \tag{5.1}
\]
The first identity follows by differentiating the top Rayleigh quotient; the
second follows by differentiation.  Convexity gives
\(h_A+h_A''\geq0\).  If this function vanished identically on every analytic
interval, \(\sigma\) would be constant on each interval and \(K\) would be a
polygon.  Otherwise analyticity supplies a subinterval on which
\(h_A+h_A''>0\).  Formula (5.1) then gives the required regular, positively curved
arc.  Differentiability of the support function makes the exposed face a
singleton there.

It remains to exclude a polygon when the spectrum is interior.  Let
\(\lambda\) be a vertex and let a unit vector \(\upsilon\) realize
\(\upsilon^*A\upsilon=\lambda\).  Two linearly independent supporting normals \(n_1,n_2\)
at the vertex give
\[
 \Rea(\overline{n_j}A)\upsilon
 =\Rea(\overline{n_j}\lambda)\upsilon
 \quad(j=1,2).
\]
The two independent real-linear equations imply
\(A\upsilon=\lambda\upsilon\) and
\(A^*\upsilon=\overline\lambda\upsilon\).  Thus \(\lambda\) is a
boundary eigenvalue, contradicting \(\spec(A)\subset\Int K\).
\end{proof}

\begin{proposition}[Rational collapse]\label{prop:rational}
Under the hypotheses and notation of Lemma~\ref{lem:extremizer}, let \(x,y\)
be the equality vectors from Proposition~\ref{prop:equality-data}.  Then the
function \(f\) is rational on \(\Omega\).  More precisely, if
\[
 a(z)=x^*(zI-A)^{-1}x,\qquad c(z)=y^*(zI-A)^{-1}x,
\]
then
\[
 c(z)f(z)=2a(z)                                       \tag{5.2}
\]
as an identity of meromorphic functions on \(\Omega\).  Equivalently,
\(f=2a/c\) wherever the quotient is defined, and then everywhere after
removable cancellation.
After cancellation, \(f=U/V\) for coprime polynomials with
\[
 \deg U>\deg V.                                       \tag{5.3}
\]
\end{proposition}

\begin{proof}
Choose the arc from Lemma~\ref{lem:arc}; it avoids the finite spectrum.  At an
almost-everywhere point \(\sigma\) of that arc, let \(\mathfrak n\) be its outward
normal, \(R=(\sigma I-A)^{-1}\), and
\[
 \Delta_\sigma=\Rea\bigl(\overline{\mathfrak n}(\sigma I-A)\bigr)\succeq0.
\]
Equations (3.5) and (3.9) imply
\[
 \Delta_\sigma R\bigl(y-f(\sigma)x\bigr)=0.           \tag{5.4}
\]
Indeed, squaring the norm in (3.9) and using
\(\mathcal P=\pi^{-1}R^*\Delta_\sigma R\) first gives
\(\Delta_\sigma^{1/2}R(y-f(\sigma)x)=0\), and hence (5.4).
The vector
\(r_\sigma=R(y-f(\sigma)x)\) is nonzero: \(R\) is invertible, while the
orthonormality of \(x,y\) makes \(y-f(\sigma)x\neq0\).  It is therefore a top
support eigenvector.  The exposed point is unique, so
\(r_\sigma^*Ar_\sigma=\sigma\|r_\sigma\|^2\).  Since
\((\sigma I-A)r_\sigma=y-f(\sigma)x\), taking a scalar product and then conjugating
gives
\[
 \bigl(y-f(\sigma)x\bigr)^*
 R\bigl(y-f(\sigma)x\bigr)=0.                         \tag{5.5}
\]

The resolvent commutes with \(T=f(A)\).  Relations (3.8) give, for every
\(z\notin\spec(A)\),
\[
 x^*R_z y=0,\qquad y^*R_z y=x^*R_zx=a(z).             \tag{5.6}
\]
Expanding (5.5), using \(|f(\sigma)|=1\) and (5.6), yields
\[
 2a(\sigma)-f(\sigma)c(\sigma)=0                     \tag{5.7}
\]
almost everywhere on the arc.  Both sides are continuous there, so (5.7)
holds on the whole arc.  On a collar of this arc the function \(cf-2a\) is
analytic, continuous to the arc, and zero there, because the spectrum is a
positive distance away.  Flatten the analytic arc to a line segment by a
local biholomorphic coordinate and extend the zero boundary function by zero
to the other side.  Morera's theorem makes the extension analytic, so the
identity theorem first gives \(cf-2a=0\) in an interior collar.  The punctured
domain \(\Omega\setminus\spec(A)\) is connected; applying the identity theorem
there proves \(cf-2a=0\) throughout it, hence as a meromorphic identity on
\(\Omega\).  Since \(a(z)=z^{-1}+O(z^{-2})\) at infinity, \(a\) is not the
zero rational function; the identity therefore also shows that \(c\) is not
identically zero.  This proves (5.2).

Write
\[
 a(z)=\frac{x^*\adj(zI-A)x}{\det(zI-A)},\qquad
 c(z)=\frac{y^*\adj(zI-A)x}{\det(zI-A)}.
\]
The numerator of \(a\) has degree exactly \(N-1\), with leading coefficient
\(1\).  Since \(x^*y=0\), the numerator of \(c\) has degree at most \(N-2\).
Canceling a common polynomial from numerator and denominator preserves the
strict degree difference, proving (5.3).
\end{proof}

\section{The full lemniscate is the boundary}

\begin{lemma}[The extremal state]\label{lem:state}
With the notation above, the functional
\[
 L(h)=x^*h(A)x,\qquad h\in A(K),
\]
is a unital positive state: if \(\Rea h\geq0\) on \(K\), then
\(\Rea L(h)\geq0\), and \(\|L\|=L(1)=1\).  Consequently there is a probability
measure \(\mu\) on \(K\) such that
\[
 L(h)=\int_Kh(\zeta)\,d\mu(\zeta)\qquad(h\in A(K)).    \tag{6.1}
\]
\end{lemma}

\begin{proof}
If \(\Rea h\geq0\), then \(f_t=f e^{-th}\) has supremum norm at most \(1\) for
\(t\geq0\).  Since \(f\) is extremal and \(Tx=2y\),
\begin{align*}
 0&\geq\left.\frac d{dt}\right|_{0+}\|f_t(A)x\|^2\\
  &=-2\Rea\bigl((Tx)^*T h(A)x\bigr)
   =-8\Rea\bigl(y^*h(A)y\bigr)
   =-8\Rea L(h),
\end{align*}
where the last equality is (3.11).  Thus \(L\) is positive on real parts and
\(L(1)=1\).  If \(\|h\|_K\leq1\), apply positivity to
\(1-e^{i\theta}h\) for every \(\theta\) to get \(|L(h)|\leq1\).  Hence
\(\|L\|=1\).

Extend \(L\) from the unital function space \(A(K)\) to \(C(K)\) by the
norm-preserving Hahn--Banach theorem.  A norm-one functional taking \(1\) to
\(1\) is positive: applying its norm bound to \(1+it g\), with \(g\) real and
letting \(t\to0\), first makes its value on \(g\) real; applying it to
\(0\leq g\leq1\) then gives a value in \([0,1]\).  The elementary Riesz
representation construction (define the mass of an open set by suprema of
continuous functions supported in it) gives a probability measure \(\mu\)
and (6.1).
\end{proof}

\begin{lemma}[Zeros of the transfer function]\label{lem:zeros}
Let \(f=U/V\) be the reduced rational function from
Proposition~\ref{prop:rational}.  Then every finite zero of \(U\) lies in
\(\Omega\), every zero of \(V\) lies outside \(K\), and \(f(\infty)=\infty\).
\end{lemma}

\begin{proof}
For \(z\notin K\), the function \(\zeta\mapsto(z-\zeta)^{-1}\) lies in
\(A(K)\).  Thus (6.1) gives
\[
 a(z)=\int_K\frac{d\mu(\zeta)}{z-\zeta}.              \tag{6.2}
\]
Strict separation of \(z\) from the convex set \(K\) supplies a unit complex
number \(\eta\) such that
\(\Rea(\overline\eta(z-\zeta))>0\) for every \(\zeta\in K\).
Therefore
\[
 \Rea\bigl(\eta a(z)\bigr)
 =\int_K\Rea\frac{\eta}{z-\zeta}\,d\mu(\zeta)>0
\]
because, on writing \(z-\zeta=\eta w\), the separation inequality says
\(\Rea w>0\), and then
\(\Rea(\eta/(z-\zeta))=\Rea(1/w)>0\).  In particular \(a(z)\neq0\) off \(K\).

The reduced numerator \(U\) of \(f=2a/c\) divides the numerator of \(a\), so it has
no zero off \(K\).  It has no zero on \(\partial K\), where \(|f|=1\), and
hence all its zeros lie in \(\Omega\).  A reduced pole cannot lie in
\(\Omega\), where \(f\) is holomorphic, or on \(\partial K\), where \(f\) has
a bounded continuous inner limit.  Thus all poles lie outside \(K\).
Finally (5.3) says exactly that \(f(\infty)=\infty\).
\end{proof}

\begin{proposition}[Full-level identity]\label{prop:full-level}
For the reduced rational extremizer \(f=U/V\),
\[
 \{z\in\C:|f(z)|<1\}=\Omega,\qquad
 \{z\in\C:|f(z)|=1\}=\partial K.                      \tag{6.3}
\]
Moreover \(f'\) does not vanish on \(\partial K\), and \(\partial K\) is a
smooth real-analytic strictly convex Jordan curve.
\end{proposition}

\begin{proof}
The domain \(\Omega\) is a component of
\(\mathcal O:=\{|f|<1\}\), because \(f\) is strictly bounded by \(1\) inside
and has modulus \(1\) on the boundary.  Every component \(\mathcal O_0\) of
\(\mathcal O\) is bounded,
since \(f(\infty)=\infty\), and it contains a zero of \(f\).  Indeed, if it
contained no zero, \(1/f\) would be holomorphic on a neighborhood of
\(\overline{\mathcal O_0}\), have modulus \(1\) on its boundary, and have modulus
strictly larger than \(1\) inside.  Its maximum would then occur in the
interior, contrary to the maximum-modulus principle.  Lemma~\ref{lem:zeros}
puts every zero in \(\Omega\), so no component other than \(\Omega\) exists.
This proves the first identity.

The boundary of \(\mathcal O\) is contained in \(\{|f|=1\}\).  Conversely, if
\(|f(z_0)|=1\), the open mapping theorem says that every neighborhood of
\(z_0\) contains points where \(|f|<1\).  Thus \(z_0\in\partial\mathcal O\), proving
the second identity.

If \(f'\) vanished to order \(k-1\geq1\) at a level point, then a local
analytic logarithm of \(f/f(z_0)\) would start with
\(c(z-z_0)^k\).  Its zero-real-part set would have \(2k\geq4\) local rays.
That is impossible for the boundary of a convex body, which is locally a
Jordan arc.  Hence \(f'\neq0\) on the level, and the implicit-function theorem
makes the boundary real analytic and smooth.

Finally suppose that the boundary contained a nontrivial segment of a real
line \(z=z_0+t\xi\), \(t\in\mathbb R\), \(\xi\neq0\).  The expression
\[
 |U(z_0+t\xi)|^2-|V(z_0+t\xi)|^2
\]
is a polynomial in the real variable \(t\).  Its vanishing on an interval
would make it vanish for every real \(t\).  Since coprime one-variable
polynomials have no common zero, the full-level identity would then put the
entire unbounded line in \(\partial K\), contradicting compactness.  Thus the
boundary has no nontrivial segment, which for a compact convex body is
exactly strict convexity.
\end{proof}

\section{The algebraic tangent argument}

Only elementary facts about plane algebraic curves are needed.  We spell them
out in the proof rather than invoke an equality classification theorem.

\begin{proposition}[A convex rational lemniscate of numerical-range type is a circle]
\label{prop:circle}
Let \(K\) be a compact convex body whose boundary is the full level set
\[
 \partial K=\{z\in\C:|U(z)/V(z)|=1\},                 \tag{7.1}
\]
where \(U,V\) are coprime polynomials, \(\deg U>\deg V\), and the level is a
smooth strictly convex Jordan curve.  Suppose that every tangent line to an
open boundary arc belongs to the projective curve
\[
 \det(sI+u\mathsf H+v\mathsf G)=0                    \tag{7.2}
\]
for two Hermitian matrices \(\mathsf H,\mathsf G\).  Then \(\partial K\) is a
circle.
\end{proposition}

\begin{proof}
Write \(z=X+iY\), \(w=X-iY\), and
\(S^\#(w)=\overline{S(\overline w)}\) for a polynomial \(S\).
For a homogenization \(S_h(w,Z)\), the notation \(S_h^\#\) means
coefficientwise conjugation, with \(w,Z\) left unchanged.
We use primal coordinates \([X:Y:Z]\); a line
\(sZ+uX+vY=0\) has dual coordinates \([s:u:v]\).  Thus, in vector
incidence formulas below, put
\(\mathbf z=(Z,X,Y)^{\mathsf T}\) and
\(\boldsymbol\ell=(s,u,v)^{\mathsf T}\).
Let \(d=\deg U\), \(e=\deg V<d\), and homogenize the real lemniscate polynomial
to degree \(2d\):
\[
 \mathscr F(X,Y,Z)=
 U_h(X+iY,Z)U_h^\#(X-iY,Z)
 -Z^{2(d-e)}V_h(X+iY,Z)V_h^\#(X-iY,Z).               \tag{7.3}
\]
Here \(U_h,V_h\) are the usual homogenizations to their own degrees.
The smooth connected curve \(\partial K\) lies in one complex irreducible
factor \(\mathscr C\) of \(\mathscr F=0\).  Indeed, locally only one factor can
vanish at a regular point; if two did, the gradient of their product would
vanish.  Connectedness keeps the same factor along the whole boundary.  Since
the boundary consists of real regular points, conjugation preserves this
factor, so its equation can be chosen real.

Every real affine point of \(\mathscr C\) is a real zero of (7.3), hence is on
the full level (7.1); coprimality prevents \(U,V\) from vanishing there
simultaneously.  Thus
\[
 \mathscr C(\mathbb R)=\partial K
\]
in the affine plane.  There are no additional real points at infinity,
because at \(Z=0\)
\[
 \mathscr F(X,Y,0)=|\operatorname{lc}(U)|^2
                    (X^2+Y^2)^d,                    \tag{7.4}
\]
which has no real projective zero.  Therefore the real projective locus of
\(\mathscr C\) is exactly the one oval \(\partial K\).

Let \(\mathscr C^*\) be the projective dual curve, the Zariski closure of the
tangent lines to \(\mathscr C\), and let
\(\mathscr G(s,u,v)\) be its irreducible homogeneous equation.  Tangency on an
open real arc and (7.2) imply
\[
 \mathscr G(s,u,v)\ \hbox{divides}\ 
 \det(sI+u\mathsf H+v\mathsf G).                     \tag{7.5}
\]
For clarity, the divisibility uses only this elementary observation: the
restriction of the determinant to a local analytic parametrization of the
irreducible dual curve vanishes on an interval, hence identically; division
with remainder in one coordinate then makes the irreducible defining
polynomial a factor.

Set \(\mathbf e=[1:0:0]\).  The determinant in (7.5) equals \(1\) at
\(\mathbf e\), so \(\mathscr G(\mathbf e)\neq0\).  For real \(u,v\), all roots
in \(s\) of the determinant are real, because they are the negatives of the
eigenvalues of the Hermitian matrix \(u\mathsf H+v\mathsf G\).  Hence every root of
\(\mathscr G(s,u,v)\) is real.  In other words, \(\mathscr G\) is hyperbolic
with respect to \(\mathbf e\).

Let \(m=\deg\mathscr G\).  Choose a generic real direction \([u:v]\).  Since
\(\mathscr G(\mathbf e)\neq0\), the coefficient of \(s^m\) in
\(\mathscr G(s,u,v)\) is nonzero.  Here are the finite exceptional sets being
avoided.  On the normalization of the non-linear curve \(\mathscr C\), the
Gauss map to \(\mathscr C^*\) is nonconstant; in characteristic zero its
derivative is not identically zero, and hence has only finitely many zeros.
The singular locus of the irreducible plane curve \(\mathscr C^*\) is also
finite.  Avoid the finitely many directions obtained by projecting those
critical images and singular points to \([u:v]\), and also avoid the zeros of
the discriminant of \(\mathscr G(s,u,v)\).  This discriminant is not
identically zero: irreducibility and characteristic zero make
\(\mathscr G\), viewed over \(\mathbb C(u/v)\), separable in \(s\).  Thus the
specialization has
\(m\) distinct roots, and
hyperbolicity makes all corresponding dual points
\(\ell_j=[s_j:u:v]\) real and smooth.

Each such \(\ell_j\) has a unique generic contact point on \(\mathscr C\), and
that point is real.  Here is a direct verification of both assertions.  If
\(\mathbf z(t)\) is a local homogeneous parametrization of \(\mathscr C\) in
the coordinate order \((Z,X,Y)\), its tangent line is
\(\boldsymbol\ell(t)=\mathbf z(t)\times\mathbf z'(t)\).  Both
\(\boldsymbol\ell(t)\) and \(\boldsymbol\ell'(t)\) annihilate
\(\mathbf z(t)\).  Hence the tangent line to \(\mathscr C^*\) at
\(\boldsymbol\ell(t)\), viewed dually, recovers \(\mathbf z(t)\).  Our roots
avoid the critical images.  If a smooth dual root had two regular contacts,
its unique tangent line would recover both primal points, so they would be
the same.  Thus the contact is unique.  Equivalently, in the standard primal
coordinate order,
it is the real projective point
\[
 [X:Y:Z]=[\mathscr G_u(\ell_j):\mathscr G_v(\ell_j):
   \mathscr G_s(\ell_j)],
\]
so it is real.

All \(m\) contacts therefore lie on
\(\mathscr C(\mathbb R)=\partial K\).  A smooth strictly convex oval has
exactly two tangent lines with a prescribed unoriented normal \([u:v]\): the
support line at the maximum and the one at the minimum of \(uX+vY\).
The preceding count shows that there are at most two roots; conversely these
two support lines belong to \(\mathscr C^*\), so both are roots.  Consequently
\[
 m=2.                                                 \tag{7.6}
\]
Thus \(\mathscr C^*\) is an irreducible conic and is nonsingular.  Writing its
equation as
\(\boldsymbol\ell^{\mathsf T}\Sigma\boldsymbol\ell=0\), with an invertible
symmetric \(3\times3\) matrix \(\Sigma\), its tangent at
\(\boldsymbol\ell\) corresponds to the primal point
\(\Sigma\boldsymbol\ell\).  The envelope of those tangents is
\(\mathbf z^{\mathsf T}\Sigma^{-1}\mathbf z=0\), also a conic.  Since it
contains a Zariski-dense
arc of \(\mathscr C\), it is \(\mathscr C\).

Finally, every point of this conic at infinity must, by divisibility in
(7.3), be one of
\[
 [1:i:0],\qquad[1:-i:0].
\]
A real irreducible conic has two points at infinity counted with
multiplicity, and conjugation exchanges these two circular points.  It
therefore contains both.  Its real homogeneous equation has the form
\[
 \alpha(X^2+Y^2)+\beta XZ+\gamma YZ+\delta Z^2=0,
 \qquad \alpha\neq0.
\]
Completing squares shows that its nonempty compact real oval is a circle.
\end{proof}

\section{Proof of the main theorem}

\begin{proof}[Proof of Theorem~\ref{thm:main}]
We use induction on \(N\).  For \(N=1\), and more generally whenever the
numerical range has empty interior, Lemma~\ref{lem:foundation} gives
\(\psi(A)=1\), so the hypothesis never occurs.

Assume the result below size \(N\), and let \(A\in M_N(\C)\) satisfy
\(\psi(A)=2\).  Put \(K=W(A)\).  It has nonempty interior.  By
Lemma~\ref{lem:induction}, either \(K\) is already a nondegenerate disk or
\(\spec(A)\subset\Omega=\Int K\).  Consider the second case.

Lemma~\ref{lem:extremizer} supplies a finite-Blaschke extremizer
\(f\in A(K)\) satisfying (4.6); thus no polynomial attainment has been
assumed.  Proposition~\ref{prop:equality-data} supplies \(x,y\) and the exact
boundary-kernel identity.  Since an interior-spectrum numerical range cannot
be a polygon, Lemma~\ref{lem:arc} supplies a curved analytic boundary arc.
On that arc, Proposition~\ref{prop:rational} proves
\[
 f(z)=\frac{U(z)}{V(z)},\qquad \deg U>\deg V,
\]
after cancellation of the meromorphic transfer identity (5.2).

The variation \(fe^{-th}\) makes \(x^*h(A)x\) a positive state
(Lemma~\ref{lem:state}).  Its Cauchy transform has no zero outside the convex
set \(K\), so every zero of the reduced numerator \(U\) lies in \(\Omega\),
while the poles lie outside \(K\).  The degree inequality makes \(f\) blow up
at infinity.  Proposition~\ref{prop:full-level} then proves the global, not
merely local, identities
\[
 \{|f|<1\}=\Omega,\qquad \{|f|=1\}=\partial K,
\]
and makes this boundary a smooth strictly convex rational lemniscate.

For \(H_A=\Rea A\) and \(G_A=\Ima A\), every supporting tangent line
\(s+uX+vY=0\) to \(W(A)\) satisfies
\[
 \det(sI+uH_A+vG_A)=0:
\]
the support matrix is semidefinite and has a unit vector attaining the
boundary point in its kernel.  All hypotheses of
Proposition~\ref{prop:circle} now hold.  It follows that \(\partial K\) is a
circle.  Since \(K\) is compact and convex, it is the filled closed disk
bounded by that circle, not merely the boundary circle itself.

The disk has positive radius because \(K\) has nonempty planar interior.
Writing its center and radius as \(\gamma\) and \(\rho\) gives
\[
 W(A)=\{z:|z-\gamma|\leq\rho\},\qquad \rho>0,
\]
which completes the induction and the proof.
\end{proof}

\section*{Context and source audit}

The proof above is logically self-contained.  For context, the double-layer
construction is historically associated with Crouzeix--Palencia, while the
telescoping estimate in Lemma~\ref{lem:dilation} is the sharp 2026
Lorist--Schwenninger mechanism.  The finite-Blaschke extremizer reduction goes
back to finite Nevanlinna--Pick interpolation, and the determinant
\(\det(sI+u\Rea A+v\Ima A)\) is the classical Kippenhahn polynomial.  These
attributions are not used as proof substitutes.  Primary formulations
consulted during the audit include:
\begin{itemize}
\item M.~Crouzeix, \emph{Bounds for analytical functions of matrices},
      Integral Equations Operator Theory 48 (2004), 461--477,
      \url{https://doi.org/10.1007/s00020-002-1188-6};
\item M.~Crouzeix and C.~Palencia, \emph{The numerical range is a
      \((1+\sqrt2)\)-spectral set}, SIAM J. Matrix Anal. Appl. 38 (2017),
      \url{https://arxiv.org/abs/1702.00668};
\item E.~Lorist and F.~L. Schwenninger, \emph{A solution to Crouzeix's
      conjecture}, \url{https://arxiv.org/abs/2608.03841};
\item B.~Malman, J.~Mashreghi, R.~O'Loughlin and T.~Ransford,
      \emph{On the Crouzeix ratio for \(N\times N\) matrices},
      \url{https://arxiv.org/abs/2409.14127};
\item S.~Orevkov and F.~Pakovich, \emph{On intersection of lemniscates of
      rational functions}, \url{https://arxiv.org/abs/2309.04983}.
\end{itemize}

\end{document}
